Many debiasing papers assume that if we can find a stereotype direction inside a model, we can remove stereotype behavior by editing that direction. We test that assumption directly by comparing matched \emph{write} and \emph{erase} edits. In our completion-format setting, write beats erase in 2 of 3 core models; in a second setting where the edited signal is at the decision token, write beats erase in all 3 core models. Results are mixed in 4 extra models (1 of 4), transfer across benchmarks is weak at current sample sizes, and one cross-source test flips direction (Llama StereoSet$\rightarrow$CrowS: $-0.350$, $q=0.0208$). The main takeaway is simple: for linear direction edits, finding a direction gives useful control, but it does not by itself prove that debiasing worked.
